% GENERATED FILE. Edit canonical lessons, then run npm run generate:books.
% canonical-source-hash: fnv1a64:a7bb087f43117c00
% canonical-lessons: UR-C91-dopahar, UR-C91-savera, UR-C91-adhirat, UR-C91-parson, UR-C91-itvar

\chapter{Midday, Early morning, Midnight, Day after tomorrow, Sunday}
\label{ch:ur-midday-early-morning-midnight-day-after-tomorrow-sunday}
\hlchaptermodality{91}

\begin{chapteropening}
\textbf{By the end of this chapter:} \emph{I can say midday, early morning, midnight, the day after tomorrow and Sunday.}

Complete the last lesson of chapter 91: I can say midday, early morning, midnight, the day after tomorrow and Sunday.
\end{chapteropening}

\section[\texorpdfstring{\ur{دوپہر} (dopahar)}{دوپہر (dopahar)}]{\ur{دوپہر} (dopahar) — midday, afternoon}

\label{lesson:UR-C91-dopahar}



\begin{quote}
Before the new one: say the Urdu for the one, then the Urdu for here.
\end{quote}

\subsection*{You'll want to know: \ur{دوپہر}}
\textbf{\ur{دوپہر}} — \emph{dopahar} — \textquotedblleft{}midday, afternoon\textquotedblright{}.

\begin{cousinweb}
\textbf{\ur{دوپہر}} (\emph{dopahar}) is \textquotedblleft{}two watches\textquotedblright{}: \emph{do}, \textquotedblleft{}two\textquotedblright{} (CDIAL 6648), and \emph{pahar}, a watch of about three hours, counted from sunrise. CDIAL lists \emph{pahar} under Sanskrit \emph{prahara}, \textquotedblleft{}a watch\textquotedblright{}, and asks whether that word is a Sanskritization of Persian \emph{pahra}, \textquotedblleft{}watch\textquotedblright{} (CDIAL 8900).
\end{cousinweb}

\begin{grammarlens}[title={What you've built}]
One more word for this chapter.
\end{grammarlens}

\subsection*{Your turn}
\begin{itemize}
\raggedright
  \item \emph{Say it:} \emph{dopahar}
  \item \emph{Say it:} \emph{dopahar}, once more
  \item \emph{Say it:} say \emph{dopahar}
  \item \emph{Recall:} say the Urdu for the one, then the Urdu for here, then say \emph{dopahar} again
\end{itemize}

\begin{tcolorbox}[breakable,colback=teal!4,colframe=teal!35!black,title={Before you move on}]
What does \ur{دوپہر} mean? (\textquotedblleft{}Midday, afternoon\textquotedblright{}.) Say it once more.
\end{tcolorbox}
\section[\texorpdfstring{\ur{سویرا} (saverā)}{سویرا (saverā)}]{\ur{سویرا} (saverā) — early morning}

\label{lesson:UR-C91-savera}



\begin{quote}
Before the new one: say the Urdu for here, then the Urdu for midday.
\end{quote}

\subsection*{You'll want to know: \ur{سویرا}}
\textbf{\ur{سویرا}} — \emph{saverā} — \textquotedblleft{}early morning\textquotedblright{}.

\begin{cousinweb}
\textbf{\ur{سویرا}} (\emph{saverā}) is a tadbhava, a word inherited and reshaped by sound change. It comes from \emph{*savela}, \textquotedblleft{}in time, early\textquotedblright{}, a form CDIAL reconstructs from its descendants (CDIAL 13291). It is built from \emph{sa-}, \textquotedblleft{}with\textquotedblright{}, and \emph{velā}, \textquotedblleft{}time\textquotedblright{}.
\end{cousinweb}

\begin{grammarlens}[title={What you've built}]
One more word for this chapter.
\end{grammarlens}

\subsection*{Your turn}
\begin{itemize}
\raggedright
  \item \emph{Say it:} \emph{saverā}
  \item \emph{Say it:} \emph{saverā}, once more
  \item \emph{Say it:} say \emph{saverā}
  \item \emph{Recall:} say the Urdu for here, then the Urdu for midday, then say \emph{saverā} again
\end{itemize}

\begin{tcolorbox}[breakable,colback=teal!4,colframe=teal!35!black,title={Before you move on}]
What does \ur{سویرا} mean? (\textquotedblleft{}Early morning\textquotedblright{}.) Say it once more.
\end{tcolorbox}
\section[\texorpdfstring{\ur{آدھی} \ur{رات} (ādhī rāt)}{آدھی رات (ādhī rāt)}]{\ur{آدھی} \ur{رات} (ādhī rāt) — midnight}

\label{lesson:UR-C91-adhirat}



\begin{quote}
Before the new one: say the Urdu for midday, then the Urdu for early morning.
\end{quote}

\subsection*{You'll want to know: \ur{آدھی} \ur{رات}}
\textbf{\ur{آدھی} \ur{رات}} — \emph{ādhī rāt} — \textquotedblleft{}midnight\textquotedblright{}.

\begin{cousinweb}
\textbf{\ur{آدھی} \ur{رات}} (\emph{ādhī rāt}) is \textquotedblleft{}half night\textquotedblright{}. \emph{ādhī}, \textquotedblleft{}half\textquotedblright{}, comes from Sanskrit \emph{ardha}, \textquotedblleft{}half\textquotedblright{} (CDIAL 644), and \emph{rāt}, \textquotedblleft{}night\textquotedblright{}, from \emph{rātrī}, \textquotedblleft{}night\textquotedblright{} (CDIAL 10702).
\end{cousinweb}

\begin{grammarlens}[title={What you've built}]
One more word for this chapter.
\end{grammarlens}

\subsection*{Your turn}
\begin{itemize}
\raggedright
  \item \emph{Say it:} \emph{ādhī rāt}
  \item \emph{Say it:} \emph{ādhī rāt}, once more
  \item \emph{Say it:} say \emph{ādhī rāt}
  \item \emph{Recall:} say the Urdu for midday, then the Urdu for early morning, then say \emph{ādhī rāt} again
\end{itemize}

\begin{tcolorbox}[breakable,colback=teal!4,colframe=teal!35!black,title={Before you move on}]
What does \ur{آدھی} \ur{رات} mean? (\textquotedblleft{}Midnight\textquotedblright{}.) Say it once more.
\end{tcolorbox}
\section[\texorpdfstring{\ur{پرسوں} (parsoṅ)}{پرسوں (parsoṅ)}]{\ur{پرسوں} (parsoṅ) — the day after tomorrow; the day before yesterday}

\label{lesson:UR-C91-parson}



\begin{quote}
Before the new one: say the Urdu for early morning, then the Urdu for midnight.
\end{quote}

\subsection*{You'll want to know: \ur{پرسوں}}
\textbf{\ur{پرسوں}} — \emph{parsoṅ} — \textquotedblleft{}the day after tomorrow; the day before yesterday\textquotedblright{}.

\begin{cousinweb}
\textbf{\ur{پرسوں}} (\emph{parsoṅ}) is a tadbhava, a word inherited from Sanskrit and reshaped by sound change. It comes from \emph{paraśvas}, \textquotedblleft{}the day after tomorrow\textquotedblright{} (CDIAL 7799j).
\end{cousinweb}

\begin{grammarlens}[title={What you've built}]
One more word for this chapter.
\end{grammarlens}

\subsection*{Your turn}
\begin{itemize}
\raggedright
  \item \emph{Say it:} \emph{parsoṅ}
  \item \emph{Say it:} \emph{parsoṅ}, once more
  \item \emph{Say it:} say \emph{parsoṅ}
  \item \emph{Recall:} say the Urdu for early morning, then the Urdu for midnight, then say \emph{parsoṅ} again
\end{itemize}

\begin{tcolorbox}[breakable,colback=teal!4,colframe=teal!35!black,title={Before you move on}]
What does \ur{پرسوں} mean? (\textquotedblleft{}The day after tomorrow; the day before yesterday\textquotedblright{}.) Say it once more.
\end{tcolorbox}
\section[\texorpdfstring{\ur{اتوار} (itvār)}{اتوار (itvār)}]{\ur{اتوار} (itvār) — Sunday}

\label{lesson:UR-C91-itvar}



\begin{quote}
Before the new one: say the Urdu for midnight, then the Urdu for the day after tomorrow; the day before yesterday.
\end{quote}

\subsection*{You'll want to know: \ur{اتوار}}
\textbf{\ur{اتوار}} — \emph{itvār} — \textquotedblleft{}Sunday\textquotedblright{}.

\begin{cousinweb}
\textbf{\ur{اتوار}} (\emph{itvār}) is usually traced to \emph{*ādityavāra}, \textquotedblleft{}Sunday\textquotedblright{}, the day of \emph{āditya}, the sun, with \emph{vāra}, \textquotedblleft{}day of the week\textquotedblright{}. CDIAL reconstructs that word from Prakrit and Sindhi forms and lists no Hindi or Urdu form (CDIAL 1154).
\end{cousinweb}

\begin{grammarlens}[title={What you've built}]
One more word for this chapter.
\end{grammarlens}

\subsection*{Your turn}
\begin{itemize}
\raggedright
  \item \emph{Say it:} \emph{itvār}
  \item \emph{Say it:} \emph{itvār}, once more
  \item \emph{Say it:} say \emph{itvār}
  \item \emph{Recall:} say the Urdu for midnight, then the Urdu for the day after tomorrow; the day before yesterday, then say \emph{itvār} again
\end{itemize}

\begin{tcolorbox}[breakable,colback=teal!4,colframe=teal!35!black,title={Before you move on}]
What does \ur{اتوار} mean? (\textquotedblleft{}Sunday\textquotedblright{}.) Say it once more.
\end{tcolorbox}
